\documentclass[10pt,a4paper]{amsart}
\usepackage[margin=19mm]{geometry}
\usepackage{amsmath,amssymb,mathtools}
\usepackage[scr=rsfs,cal=euler]{mathalfa}
\usepackage{xfrac}
\usepackage[unicode,hidelinks,bookmarks=false]{hyperref}
\newcommand{\ie}{i.e.\ }
\newcommand{\cf}{cf.\ }
\newcommand{\resp}{respectively\ }
\newcommand{\Subsection}{\subsection}
\providecommand{\nobreakdash}{\nobreak\hbox{-}}
\setlength{\emergencystretch}{2em}
\multlinegap0pt
\usepackage{amssymb}
\usepackage{enumerate}
\usepackage{stackengine}
\usepackage{xcolor}
\usepackage{caption}
\usepackage{tikz}
\usepackage{mathtools}
\newtheorem{lemma}{Lemma}
\def\f{\mathbf{f}}
\def\p{\mathbf{p}}
\def\q{\mathbf{q}}
\title{Adaptive dynamics of alternating prisoner's dilemma with memory $N$}
\author{Nataliya A. Balabanova, Hong Duong, Christian Hilbe}
\date{Source: arXiv:2601.18671v1 (CC BY 4.0)}
\begin{document}
\maketitle

\noindent\emph{Source and licence.} Adaptive dynamics of alternating prisoner's dilemma with memory $N$. Version arXiv:2601.18671v1 (CC BY 4.0), distributed by arXiv under the Creative Commons Attribution 4.0 International licence, http://creativecommons.org/licenses/by/4.0/, as recorded in the arXiv licence metadata for this exact version and shown on the abstract page https://arxiv.org/abs/2601.18671v1. Full licence notice: \texttt{licenses/arxiv-2601.18671-CC-BY-4.0.txt}.\par\smallskip

\noindent\textit{Editorial scope.} Counted unit: the article's lemma defined (source0.tex lines 230--233). The article's complete original proof of it is delivered with it (source0.tex lines 234--279, one \texttt{proof} environment of 46 lines). No mathematics is written, repaired or re-derived: the statement and the proof are the article's own lines, in the article's own notation, and no retained line is substituted or re-flowed.  No further source-local context is needed: every cross-reference of every retained line resolves inside the two delivered spans.  Nothing was omitted from any retained span, so the package carries no omission record.  No retained line contains a citation, so the wrapper carries no bibliography and no bibliography entry.  No retained line contains a figure environment, an image-inclusion command or drawing code, so no image is delivered and the package declares no asset dependency.  Editorial formatting only, in addition: an AMS \texttt{amsart} preamble that restates the article's own declarations and macros used by the retained lines, namely the environments \texttt{lemma} on separate counters, together with the standard packages those lines need.  The retained environments print in this document as Lemma 1 in order of appearance, whereas the article prints the same items with the numbering of its own full document.  The complete native source of the article as distributed in the arXiv e-print for this exact version is bundled unchanged as \texttt{sources/193405/source0.tex}.
\begin{lemma}
\label{st: payoff vector correctly defined}
    The vector $\f_N$ is well-defined, i.e. the winnings of the two players are equal for identical game histories. 
\end{lemma}

\begin{proof}
    We start with the toy case $N=1$. Then the two strategy vectors and the payoff vectors of the winnings of $\p$ and $\q$ respectively  are
    \[
    \p:=\begin{pmatrix}
        p_{CC}\\
        p_{CD}\\
        p_{DC}\\
        p_{DD}
    \end{pmatrix}, \ \ \q:=\begin{pmatrix}
        q_{CC}\\
        q_{CD}\\
        q_{DC}\\
        q_{DD}
    \end{pmatrix}, \ \ \f_{\p}:=\begin{pmatrix}
        f_{CC}^{\p}\\
        f_{CD}^{\p}\\
        f_{DC}^{\p}\\
        f_{DD}^{\p}
    \end{pmatrix}, \ \ \f_{\q}:=\begin{pmatrix}
        f_{CC}^{\q}\\
        f_{CD}^{\q}\\
        f_{DC}^{\q}\\
        f_{DD}^{\q}
    \end{pmatrix}
    \]

    Consider $f_{\p}$ first. The entry $f_{CC}^{\p}$ denotes the winnings of $\p$ after one complete round, where $\p$ and $\q$ both chose to cooperate. On the other hand, the element $f^{\q}_{CC}$ encodes the winnings of $\q$ after $\q$ decided to cooperate in the previous round and $\p$ cooperated in the current round. Using the notation from REF, we will compute both vectors:
    \[
    f_{\p} =\begin{pmatrix}
        a + b\\
        a + d\\
        c + b\\
        d + c
    \end{pmatrix},    f_{\q} =\begin{pmatrix}
      a + b\\
       a + d\\
       c+b\\
       c+d
    \end{pmatrix}
    \] 
    and $f_{\p} = \f_{\q}$. 

    For general $N$ (as, actually, for $N=1$) the proof is very straightforward. Consider two elements of $\p$ and $\q$ with the same indices: $p_{i_1i_2\ldots i_{2N}}$ and $q_{i_1i_2\ldots i_{2N}}$.  the index of $\p$ encodes first action of $\p$, then action of $\q$, then action of $\p$ and so on.  But in the index of $\q$ their action comes first, then goes the action of $\p$, and, again, so on. 

    Due to how the payoffs per choice are defined and their inherent symmetries, the payoffs after $2N$ rounds are the same. 
\end{proof}

\end{document}
